\documentclass[sigconf, final, screen]{acmart}
\AtBeginDocument{%
  \providecommand\BibTeX{{%
    \normalfont B\kern-0.5em{\scshape i\kern-0.25em b}\kern-0.8em\TeX}}}
\usepackage{times}
\usepackage{epsfig}
\usepackage{graphicx}
\usepackage{amsmath}
\let\Bbbk\relax 
\usepackage{amssymb}
\usepackage{algorithm, algpseudocode}
\usepackage{multirow}
\usepackage{caption}
\usepackage{comment}
\usepackage{amsmath}
\usepackage{amssymb}
\newcommand{\tabincell}[2]{\begin{tabular}{@{}#1@{}}#2\end{tabular}}
\usepackage{pifont}
\usepackage{color}
\usepackage{makecell}
\usepackage{booktabs}
\usepackage{tabularx}
\newcommand{\cmark}{\ding{51}}%
\newcommand{\xmark}{\ding{55}}%
\def\httilde{\mbox{\tt\raisebox{-.5ex}{\symbol{126}}}}
\newlength\savewidth\newcommand\shline{\noalign{\global\savewidth\arrayrulewidth
  \global\arrayrulewidth 1pt}\hline\noalign{\global\arrayrulewidth\savewidth}}
\newcommand{\zznote}[1]{\textcolor{blue}{ZZ: #1}}
\def\eg{\emph{e.g.}} 
\def\ie{\emph{i.e.}} 
\def\etal{\emph{et~al.}} 

% ====================================== %

\begin{document}

\settopmatter{printacmref=false}
\renewcommand\footnotetextcopyrightpermission[1]{}
\setcopyright{none}
\pagestyle{plain}

\title{ CISC7202 Practical Machine Learning \\ Experiment Report-3 }

\author{  
Yongchao Lai } 
\affiliation{
  \institution{ICI, University of Macau, Macau, China}
  Student ID: MC564057}
\email{yongchao.lai@connect.um.edu.mo}

\begin{abstract}

\end{abstract}
 
\maketitle

\section{Introduction}

\section{Experimental Setup}

\subsection{Environment}
All experiments were conducted on a cloud server with an \textbf{NVIDIA RTX 3090 GPU (24GB VRAM)}. The software stack used \textbf{Python 3.9}, \textbf{PyTorch 1.13.1}, and \textbf{CUDA 11.6}. The compiler toolchain consisted of \textbf{GCC 11.5} and \textbf{G++ 11.5}.

\subsection{Baseline}

% The baseline experiment uses the text prompt ``a lion, animated movie character, high detail 3d model.'' All other settings follow the default configuration of the project unless otherwise specified.

% The key parameters are as follows: the number of training iterations is set to 6000, the learning rate is $1\times10^{-5}$, the rendering resolution for NeRF during training is $64\times64$, and Stable Diffusion 2.1 is used as the guidance model. This baseline serves as the reference configuration for subsequent controlled experiments.

The baseline is defined using the prompt ``a lion, animated movie character, high detail 3d model,'' while all other settings adopt the default configuration provided in the project. The main training parameters are listed in Table~\ref{tab:baseline_config}. Baseline experiment configuration used as the reference setup for all controlled experiments.

\begin{table}[h]
\caption{Baseline experiment configuration}
\label{tab:baseline_config}
\centering
\begin{tabularx}{\linewidth}{lX}
\toprule
\textbf{Parameter} & \textbf{Value} \\
\midrule
Text Prompt & a lion, animated movie character, high detail 3d model \\
Training Iterations & 6000 \\
Learning Rate & $1 \times 10^{-5}$ \\
NeRF Rendering Resolution & $64 \times 64$ \\
Guidance Model & Stable Diffusion 2.1 \\
\bottomrule
\end{tabularx}
\end{table}

\subsection{DMTet Fine-tuning}
The DMTet experiment was designed as a two-stage training procedure in order to compare a standard baseline run with its mesh-based fine-tuning counterpart. Two configurations were considered in this section: \textit{baseline} and \textit{baseline + DMTet}. In the baseline setting, the model was trained with the standard NeRF-based optimization under the default project configuration. In the baseline + DMTet setting, a second-stage DMTet fine-tuning process was performed by initializing from the checkpoint produced by the first stage and then enabling \texttt{--dmtet}.

For this group of experiments, the text prompt was fixed as ``a capybara wearing a top hat, low poly.'' The first stage served as the coarse text-to-3D generation stage and was trained for 10000 iterations. The second stage used DMTet fine-tuning for an additional 5000 iterations, initialized from the checkpoint of the first-stage baseline run. Unless otherwise stated, the remaining settings followed the default implementation configuration.

Compared with the baseline stage, DMTet fine-tuning changed the optimization target from volumetric NeRF training to mesh-oriented refinement. Under the default DMTet setting in this project, the training resolution was increased from $64\times64$ to $512\times512$, which provided a higher-resolution supervision signal during the fine-tuning stage.

\begin{table}[h]
\caption{Configuration of the DMTet fine-tuning experiment}
\label{tab:dmtet_config}
\centering
\begin{tabularx}{\linewidth}{lXX}
\toprule
\textbf{Parameter} & \textbf{Baseline} & \textbf{Baseline + DMTet} \\
\midrule
Text Prompt & a capybara wearing a top hat, low poly & a capybara wearing a top hat, low poly \\
Stage Description & First-stage NeRF training & Second-stage DMTet fine-tuning \\
Initialization & Random / default initialization & Initialized from the baseline checkpoint \\
Training Iterations & 10000 & 5000 \\
Training Resolution & $64 \times 64$ & $512 \times 512$ \\
Guidance Model & Stable Diffusion 2.1 & Stable Diffusion 2.1 \\
\bottomrule
\end{tabularx}
\end{table}

\subsection{Prompt Variant Experiment}
The prompt-variant experiment was designed to evaluate model performance under progressively more constrained textual descriptions. Starting from a baseline prompt, additional semantic constraints were incrementally introduced, including role identity, clothing attributes, color specification, and object interaction. To maintain clarity and manageability, a reduced set of representative prompts was selected.

\begin{table}[h]
\caption{Reduced prompt variants used in the prompt complexity experiment}
\label{tab:prompt_variants}
\centering
\begin{tabularx}{\linewidth}{cX}
\toprule
ID & Prompt \\
\midrule
1 & a lion, animated movie character, high detail 3d mode \\
2 & a lion mayor, animated movie character, wearing a suit and tie, high detail 3d mode \\
3 & a confident lion mayor, animated movie character, wearing a blue suit and red tie, holding a donut and a cup of coffee, high detail 3d mode \\
\bottomrule
\end{tabularx}
\end{table}

The prompts were arranged in increasing order of semantic specificity. Prompt 1 corresponds to the baseline setting with minimal description. Prompt 2 introduces role identity and basic clothing attributes. Prompt 3 further adds explicit color information as well as object interaction, resulting in the most constrained description among the selected variants.

\subsection{Hyperparameter Study}
This study investigates the impact of two key training factors: learning rate and training resolution. The resolution setting is controlled by jointly adjusting \texttt{--w} and \texttt{--h}, which denote the render width and height during the NeRF training stage. In MTN, the default resolution is $64\times64$.

To ensure a controlled comparison, \texttt{--w} and \texttt{--h} were always varied together so that the rendered images remained square. All other training settings were kept fixed, including the baseline prompt, 6000 iterations, Adan optimizer, Stable Diffusion 2.1 guidance, batch size 1, and random seed 3407.

\paragraph{Learning Rate Study}
The learning rate ablation evaluates two configurations: the baseline setting of $1\times10^{-5}$ and a higher value of $5\times10^{-4}$, while keeping the resolution fixed at $64\times64$.

\begin{table}[h]
\caption{Learning rate variants}
\label{tab:lr_variants}
\centering
\begin{tabularx}{\linewidth}{c c X}
\toprule
ID & Learning Rate & Description \\
\midrule
L1 & $1\times10^{-5}$ & Baseline configuration. \\
L2 & $5\times10^{-4}$ & Higher learning rate. \\
\bottomrule
\end{tabularx}
\end{table}

\paragraph{Resolution Study}
The resolution ablation evaluates three configurations by varying \texttt{--w} and \texttt{--h} together: $32\times32$, $64\times64$ (baseline), and $128\times128$. The learning rate is fixed at the baseline value of $1\times10^{-5}$.

\begin{table}[h]
\caption{\texttt{--w}/\texttt{--h} variants used in the training-resolution ablation}
\label{tab:wh_variants}
\centering
\begin{tabularx}{\linewidth}{c c X}
\toprule
ID & Resolution & Description \\
\midrule
R1 & $32\times32$ & Lower-resolution configuration. \\
R2 & $64\times64$ & Baseline configuration. \\
R3 & $128\times128$ & Higher-resolution configuration within the GPU budget. \\
\bottomrule
\end{tabularx}
\end{table}
 
The baseline configuration is defined as a learning rate of $1\times10^{-5}$ with a resolution of $64\times64$. Each ablation study varies only one factor at a time relative to this baseline.

% \subsection{IF Guidance}
% In the IF-guidance experiment, Stable Diffusion was replaced by DeepFloyd-IF as the text guidance model. All other settings were kept aligned with the baseline configuration as closely as possible. Under this setting, IF provided the diffusion-based supervision signal during 3D optimization.

\section{Acknowledgments}
This project employed AI tools,\href{https://chatgpt.com}{ChatGPT}, to assist with information retrieval, text polishing and refinement, and project code improvements. All AI‑generated content was carefully reviewed, validated, and integrated by me, and has been adopted as part of the final outcomes of this report.

{\small
\bibliographystyle{ACM-Reference-Format}
\bibliography{references}
}

\end{document}
